\chapter{オペアンプによる演算回路}

\section{はじめに}

【目的】

本章では演算増幅器（オペアンプ）による演算回路について理解します。

【目標】

\begin{itemize}
    \item オペアンプの原理、等価回路、イマジナリーショートについて説明できること。
    \item オペアンプを用いた各演算回路について原理および動作を説明できること。
    \item 各演算回路を SPICE シミュレータによってシミュレーションできること。
    \item 電子回路を実装および計測することでシミュレーションと実機の違いについて理解すること。
    \item 各演算回路を組み合わせることで人工ニューロンおよびMLPを電子回路で実装できることを理解し、自らの手で実装できること。
\end{itemize}

【章構成】

\begin{enumerate}
    \item オペアンプ
    \item オペアンプを用いた演算回路
    \item 実験1：LTSPICE によるシミュレーション
    \item 実験2：電子回路への実装と計測
    \item 共同実験1：電子回路によるMLPの実装
    \item レポート2
\end{enumerate}

\section{オペアンプ}

% オペアンプについての説明

% オペアンプの等価回路

% イマジナリーショート

\section{オペアンプを用いた演算回路}

\subsection{反転増幅回路}

\subsection{非反転増幅回路}

\subsection{加算回路}

\subsection{減算回路}

\subsection{加減算回路}

\subsection{理想ダイオード}

\subsection{ボルテージフォロワ}

\section{LTSPICE によるシミュレーション}

% LTSPICE のセットアップは付録へ

% 加算回路：DC成分

% 加算回路：AC 成分

% 加算回路：電源電圧によるクリッピング

\section{電子回路への実装と計測}

% 加算回路の実装

% 加算回路の動作確認

% 理想ダイオードの実装

% 理想ダイオードの動作確認

\section{電子回路によるMLPの実装}

% 実装 0-0 側

% 実装 0-1 側

% 実装 1 層目

% 動作確認

\section{レポート2}

作成した MLP 回路を用いた実験結果についてレポートにまとめよ。
レポートは穴埋めとして用意してあるため、実験データおよび考察を埋めること。